% GENERATED FILE. Edit canonical lessons, then run npm run generate:books.
% canonical-source-hash: fnv1a64:712a8db264c23ad5
% canonical-lessons: TE-C527-sadhincu, TE-C527-ranincu, TE-C527-viphalamavu, TE-C527-rajinama, TE-C527-maneyu

\chapter{Success and Failure}
\label{ch:te-success-and-failure}
\hlchaptermodality{527}

\begin{chapteropening}
\textbf{By the end of this chapter:} \emph{I can say that someone achieved, excelled, failed, resigned or quit.}

Complete the last lesson of chapter 527: I can say that someone achieved, excelled, failed, resigned or quit.
\end{chapteropening}

\section[\texorpdfstring{\te{సాధించు} (sādhiñcu)}{సాధించు (sādhiñcu)}]{\te{సాధించు} (sādhiñcu) — to achieve; to accomplish}

\label{lesson:TE-C527-sadhincu}



\begin{quote}
Before the new one: say the Telugu for to cooperate, then the Telugu for training.
\end{quote}

\subsection*{You'll want to know: \te{సాధించు}}
\textbf{\te{సాధించు}} — \emph{sādhiñcu} — \textquotedblleft{}to achieve; to accomplish\textquotedblright{}.

\begin{quote}
\te{పట్టుదలతో} \te{ఆమె} \te{తన} \te{లక్ష్యం} \te{సాధించింది}.
\end{quote}

\begin{quote}
\emph{paṭṭudalatō āme tana lakṣyaṁ sādhiñcindi.}
\end{quote}

\begin{quote}
With determination she achieved her goal.
\end{quote}

Past (I …): \te{సాధించాను}, \emph{sādhiñcānu}.

Its noun \te{సాధన} is practice.

\begin{cousinweb}
\textbf{\te{సాధించు}} (\emph{sādhiñcu}) is built on Sanskrit \emph{sādhana}, \textquotedblleft{}furthering; completion\textquotedblright{} (CDIAL 13334), with the Telugu verb ending -\te{ఇంచు}.
\end{cousinweb}

\begin{grammarlens}[title={What you've built}]
One more word for describing what you lived and felt.
\end{grammarlens}

\subsection*{Your turn}
\begin{itemize}
\raggedright
  \item \emph{Say it:} \emph{sādhiñcu}
  \item \emph{Say it:} \emph{sādhiñcu}, once more
  \item \emph{Say it:} \emph{paṭṭudalatō āme tana lakṣyaṁ sādhiñcindi.} — With determination she achieved her goal.
  \item \emph{Recall:} say the Telugu for to cooperate, then the Telugu for training, then say \emph{sādhiñcu} again
\end{itemize}

\begin{tcolorbox}[breakable,colback=teal!4,colframe=teal!35!black,title={Before you move on}]
What does \te{సాధించు} mean? (\textquotedblleft{}To achieve; to accomplish\textquotedblright{}.) Say it once more.
\end{tcolorbox}
\section[\texorpdfstring{\te{రాణించు} (rāṇiñcu)}{రాణించు (rāṇiñcu)}]{\te{రాణించు} (rāṇiñcu) — to excel; to shine}

\label{lesson:TE-C527-ranincu}



\begin{quote}
Before the new one: say the Telugu for training, then the Telugu for to achieve.
\end{quote}

\subsection*{You'll want to know: \te{రాణించు}}
\textbf{\te{రాణించు}} — \emph{rāṇiñcu} — \textquotedblleft{}to excel; to shine\textquotedblright{}.

\begin{quote}
\te{మా} \te{అమ్మాయి} \te{చదువులో} \te{బాగా} \te{రాణిస్తోంది}.
\end{quote}

\begin{quote}
\emph{mā ammāyi caduvulō bāgā rāṇistōndi.}
\end{quote}

\begin{quote}
Our daughter is doing very well at her studies.
\end{quote}

Past (I …): \te{రాణించాను}, \emph{rāṇiñcānu}.

-\te{లో} \te{రాణించు}, to excel at.

\begin{cousinweb}
\textbf{\te{రాణించు}} (\emph{rāṇiñcu}) is a Telugu verb for doing well and standing out; its origin is not settled.
\end{cousinweb}

\begin{grammarlens}[title={What you've built}]
One more word for describing what you lived and felt.
\end{grammarlens}

\subsection*{Your turn}
\begin{itemize}
\raggedright
  \item \emph{Say it:} \emph{rāṇiñcu}
  \item \emph{Say it:} \emph{rāṇiñcu}, once more
  \item \emph{Say it:} \emph{mā ammāyi caduvulō bāgā rāṇistōndi.} — Our daughter is doing very well at her studies.
  \item \emph{Recall:} say the Telugu for training, then the Telugu for to achieve, then say \emph{rāṇiñcu} again
\end{itemize}

\begin{tcolorbox}[breakable,colback=teal!4,colframe=teal!35!black,title={Before you move on}]
What does \te{రాణించు} mean? (\textquotedblleft{}To excel; to shine\textquotedblright{}.) Say it once more.
\end{tcolorbox}
\section[\texorpdfstring{\te{విఫలమవు} (viphalamavu)}{విఫలమవు (viphalamavu)}]{\te{విఫలమవు} (viphalamavu) — to fail; to be unsuccessful}

\label{lesson:TE-C527-viphalamavu}



\begin{quote}
Before the new one: say the Telugu for to achieve, then the Telugu for to excel.
\end{quote}

\subsection*{You'll want to know: \te{విఫలమవు}}
\textbf{\te{విఫలమవు}} — \emph{viphalamavu} — \textquotedblleft{}to fail; to be unsuccessful\textquotedblright{}.

\begin{quote}
\te{మొదటి} \te{ప్రయత్నంలో} \te{మేము} \te{విఫలమయ్యాము}.
\end{quote}

\begin{quote}
\emph{modaṭi prayatnaṁlō mēmu viphalamayyāmu.}
\end{quote}

\begin{quote}
We failed at the first attempt.
\end{quote}

Past (I …): \te{విఫలమయ్యాను}, \emph{viphalamayyānu}.

Formal: \te{విఫలం}, fruitless, + \te{అవు}, to become.

\begin{cousinweb}
\textbf{\te{విఫలమవు}} (\emph{viphalamavu}) is built from Sanskrit \emph{viphala}, \textquotedblleft{}fruitless\textquotedblright{} (CDIAL 11803) and the Telugu verb \te{అవు}, \textquotedblleft{}to become\textquotedblright{}.
\end{cousinweb}

\begin{grammarlens}[title={What you've built}]
One more word for describing what you lived and felt.
\end{grammarlens}

\subsection*{Your turn}
\begin{itemize}
\raggedright
  \item \emph{Say it:} \emph{viphalamavu}
  \item \emph{Say it:} \emph{viphalamavu}, once more
  \item \emph{Say it:} \emph{modaṭi prayatnaṁlō mēmu viphalamayyāmu.} — We failed at the first attempt.
  \item \emph{Recall:} say the Telugu for to achieve, then the Telugu for to excel, then say \emph{viphalamavu} again
\end{itemize}

\begin{tcolorbox}[breakable,colback=teal!4,colframe=teal!35!black,title={Before you move on}]
What does \te{విఫలమవు} mean? (\textquotedblleft{}To fail; to be unsuccessful\textquotedblright{}.) Say it once more.
\end{tcolorbox}
\section[\texorpdfstring{\te{రాజీనామా} (rājīnāmā)}{రాజీనామా (rājīnāmā)}]{\te{రాజీనామా} (rājīnāmā) — a resignation}

\label{lesson:TE-C527-rajinama}



\begin{quote}
Before the new one: say the Telugu for to excel, then the Telugu for to fail.
\end{quote}

\subsection*{You'll want to know: \te{రాజీనామా}}
\textbf{\te{రాజీనామా}} — \emph{rājīnāmā} — \textquotedblleft{}a resignation\textquotedblright{}.

\begin{quote}
\te{జీతం} \te{సరిపోక} \te{అతను} \te{రాజీనామా} \te{చేశాడు}.
\end{quote}

\begin{quote}
\emph{jītaṁ saripōka atanu rājīnāmā cēśāḍu.}
\end{quote}

\begin{quote}
His pay was not enough, so he resigned.
\end{quote}

\te{రాజీనామా} \te{చేయు}, to resign.

\begin{cousinweb}
\textbf{\te{రాజీనామా}} (\emph{rājīnāmā}) is a loan from Persian \emph{rāzīnāma}, \textquotedblleft{}a written agreement, a letter of consent\textquotedblright{}, which reached Telugu through Hindustani.
\end{cousinweb}

\begin{grammarlens}[title={What you've built}]
One more word for describing what you lived and felt.
\end{grammarlens}

\subsection*{Your turn}
\begin{itemize}
\raggedright
  \item \emph{Say it:} \emph{rājīnāmā}
  \item \emph{Say it:} \emph{rājīnāmā}, once more
  \item \emph{Say it:} \emph{jītaṁ saripōka atanu rājīnāmā cēśāḍu.} — His pay was not enough, so he resigned.
  \item \emph{Recall:} say the Telugu for to excel, then the Telugu for to fail, then say \emph{rājīnāmā} again
\end{itemize}

\begin{tcolorbox}[breakable,colback=teal!4,colframe=teal!35!black,title={Before you move on}]
What does \te{రాజీనామా} mean? (\textquotedblleft{}A resignation\textquotedblright{}.) Say it once more.
\end{tcolorbox}
\section[\texorpdfstring{\te{మానేయు} (mānēyu)}{మానేయు (mānēyu)}]{\te{మానేయు} (mānēyu) — to quit; to stop}

\label{lesson:TE-C527-maneyu}



\begin{quote}
Before the new one: say the Telugu for to fail, then the Telugu for a resignation.
\end{quote}

\subsection*{You'll want to know: \te{మానేయు}}
\textbf{\te{మానేయు}} — \emph{mānēyu} — \textquotedblleft{}to quit; to stop\textquotedblright{}.

\begin{quote}
\te{అన్నయ్య} \te{పాత} \te{ఉద్యోగం} \te{మానేశాడు}.
\end{quote}

\begin{quote}
\emph{annayya pāta udyōgaṁ mānēśāḍu.}
\end{quote}

\begin{quote}
My elder brother quit his old job.
\end{quote}

Past (I …): \te{మానేశాను}, \emph{mānēśānu}.

With a job, a habit or study: \te{చదువు} \te{మానేయు}, to drop out of school.

\begin{cousinweb}
\textbf{\te{మానేయు}} (\emph{mānēyu}) is native Dravidian: DEDR 4811 lists Telugu \emph{mānu}, \textquotedblleft{}to cease, to stop, to give up, to quit\textquotedblright{}; \emph{vēyu}, \textquotedblleft{}to put\textquotedblright{}, makes the stopping final.
\end{cousinweb}

\begin{grammarlens}[title={What you've built}]
One more word for describing what you lived and felt.
\end{grammarlens}

\subsection*{Your turn}
\begin{itemize}
\raggedright
  \item \emph{Say it:} \emph{mānēyu}
  \item \emph{Say it:} \emph{mānēyu}, once more
  \item \emph{Say it:} \emph{annayya pāta udyōgaṁ mānēśāḍu.} — My elder brother quit his old job.
  \item \emph{Recall:} say the Telugu for to fail, then the Telugu for a resignation, then say \emph{mānēyu} again
\end{itemize}

\begin{tcolorbox}[breakable,colback=teal!4,colframe=teal!35!black,title={Before you move on}]
What does \te{మానేయు} mean? (\textquotedblleft{}To quit; to stop\textquotedblright{}.) Say it once more.
\end{tcolorbox}
